% !TEX root = three_object_window.tex
\documentclass[a4paper,11pt]{ltjsarticle}
\usepackage[margin=23mm,headheight=16pt]{geometry}
\usepackage{amsmath,amssymb,amsthm,mathtools}
\usepackage[haranoaji]{luatexja-preset}
\usepackage[unicode,hidelinks]{hyperref}
\usepackage{fancyhdr}
\newtheorem{theorem}{定理}[section]
\newtheorem{lemma}[theorem]{補題}
\newtheorem{corollary}[theorem]{系}
\newtheorem{definition}[theorem]{定義}
\newtheorem{remark}[theorem]{注意}
\renewcommand{\proofname}{証明}
\numberwithin{equation}{section}
\setlength{\emergencystretch}{2em}
\pagestyle{fancy}\fancyhf{}
\fancyhead[L]{\small 三対象windowと局所比較}
\fancyhead[R]{\small 2026-11-12 / B-v42}
\fancyfoot[C]{\thepage}
\makeatletter
\def\input@path{{papers/support/}}
\AtBeginDocument{\let\originalref\ref
\renewcommand{\ref}[1]{\ifcsname r@#1\endcsname\originalref{#1}\else\texttt{\detokenize{#1}}\fi}}
\makeatother
\title{三つのstabilizationによる局所残余圏の表示\\\large 論文Bの局所比較仮定の除去}
\author{}
\date{2026年11月12日 / B-v42・証明抜粋}
\begin{document}
\maketitle
\noindent 本文外の節への参照はラベル名で表示する。原稿全体の正本は
\texttt{uedakazushi/nc-cubic}の\texttt{papers/paper\_b.tex}である。
本抜粋の組版と、論文B全文の組版・全証明の独立監査は別の検証である。
\input{papers/reconstruction/three_object_window}
\begin{thebibliography}{9}
\bibitem{LP-BS}
J. Burke and G. Stevenson,
\emph{The derived category of a graded Gorenstein ring},
arXiv:1507.00830v1 (2015).
Proposition 4.18, Lemmas 6.1--6.2, Theorem 6.4.
\url{https://arxiv.org/abs/1507.00830}.
\bibitem{AMChen}
X.-W. Chen, \emph{Module factorizations},
arXiv:2402.11613v3 (27 August 2025), Theorem 7.2。
\url{https://arxiv.org/abs/2402.11613}.
\input{papers/reconstruction/three_object_window_references}
\end{thebibliography}
\end{document}
